% ---
% title: GRPO, DAPO and λGRPO
% date: 2026/1/3
% tag: RL
% ---
This blog is a summary of GRPO, DAPO and λGRPO.

\section{GRPO}

\[
\mathcal{J}_{GRPO}(\theta) = \frac{1}{\sum_{i=1}^{G} \mid o_i \mid} \sum_{i=1}^{G} \cdot \textcolor{yellow}{\frac{\text{mean}(o)}{\mid o_i \mid}} \cdot \sum_{t=1}^{|o_i|} r_{i,t}(\theta) \widehat{A_{i,t}}
\]

\section{DAPO}

\[
\mathcal{J}_{DAPO}(\theta) = \frac{1}{\sum_{i=1}^{G} \mid o_i \mid} \sum_{i=1}^{G} \cdot \text{\textcolor{green}{1}} \cdot \sum_{t=1}^{|o_i|} r_{i,t}(\theta) \widehat{A_{i,t}}
\]

\section{\(\lambda\)-GRPO}

\[
\mathcal{J}_{\lambda-\mathit{GRP}O}(\theta) = \frac{1}{\sum_{i=1}^{G} |o_i|} \sum_{i=1}^{G} \cdot \textcolor{blue}{\mathbf{f}(o_i, \lambda)} \cdot \sum_{t=1}^{|o_i|} r_{i,t}(\theta) \widehat{A_{i,t}}
\]

\textbf{GRPO 不像传统 RLHF 那样额外训练“价值模型(value)”和“奖励模型(reward model)”，而是直接用一个规则验证器(verifier) 给答案判对/判错（或给分）来当训练信号。}

\section{这些名词分别是什么}

\textbf{Policy（策略/生成模型）}

就是你的 LLM，本质是 \(\pi_\theta(y|x)\)：给定问题 \(x\)，生成回答\(y\)。

\textbf{**Reward Model（奖励模型 RM）}**

一个单独训练的模型 \(R_\phi(x,y)\)，输入“问题+回答”，输出一个标量分数，近似“人类偏好/好坏”。

\begin{itemize}
  \item 在经典 RLHF 里，人类做成对偏好数据（\(A\) 比\(B\) 好），RM 学会打分。
\end{itemize}

\textbf{Value Model（价值模型/critic）}

用来估计“从当前生成状态继续下去，最终能拿到多大回报”，记作 \(V_\psi(\text{state})\)。

\begin{itemize}
  \item 在语言生成里，\(state\) 可以理解为“\(prompt\) + 已生成前缀”。
  \item 它的主要作用：\textbf{做基线(baseline)降低方差}，算 \(advantage\)：\(A = R - V\)
\end{itemize}

\textbf{Verifier（规则验证器）}

不是学出来的偏好模型，而是“可程序化/可验证”的判分器。

\begin{itemize}
  \item 典型：数学题用答案是否正确、是否满足格式；代码题用单元测试是否通过；推理题用规则检查。
  \item 输出通常是 0/1 或可解释的分数。
\end{itemize}

\section{以往（PPO/RLHF）怎么做？}

典型 RLHF + PPO 流程（简化版）：

\begin{enumerate}
  \item \textbf{采样}：用当前策略生成回答 \(y\)。
  \item \textbf{打分}：用奖励模型 RM 得到 \(R_\phi(x,y)\)。
  \item \textbf{价值估计}：价值模型给 \(V_\psi(x, y_{<t})\)。
  \item \textbf{算优势}：\(A_t \approx R - V_t\)（可能还有 GAE 等）。
  \item \textbf{PPO 更新}：用 \(clip\) 目标更新 \(policy\)，同时通常还会加 \textbf{KL 约束}，别偏离参考模型太远。
\end{enumerate}

问题是：

\begin{itemize}
  \item 你得训练/维护 \textbf{RM + Value} 两个额外模型（算力和工程复杂度都高）；
  \item PPO 本身也比较敏感、容易不稳定。
\end{itemize}

\section{GRPO 为什么说“去掉 value model 和 reward model”？}

\textbf{去掉 Reward Model}：

GRPO 不用“学出来的偏好打分器”，而用 \textbf{规则 verifier} 直接给 reward（例如“答案对=1，不对=0”）。

\textbf{去掉 Value Model}：

GRPO 用“\textbf{组内相对比较}”来当 \(baseline\)，替代 \(critic/value\)。做法是：

\begin{itemize}
  \item 对同一个问题 \(x\)，一次生成一组候选回答\( \{y^{(1)},...,y^{(G)}\}\)
  \item verifier 分别打分 \(\{r^{(i)}\}\)
  \item 用组内均值/标准差构造相对优势（直觉：比同组平均好就正优势，比平均差就负优势），例如：
  
  \[
  \hat A^{(i)} = \frac{r^{(i)} - \text{mean}(r)}{\text{std}(r)+\epsilon}
  \]
\end{itemize}

这样就不需要训练一个 \(V_\psi\) 来做 \(baseline\) 了——\(baseline\) 直接来自“同题同组的其他回答”。

\section{一句话对比总结}

\begin{itemize}
  \item \textbf{PPO/RLHF}：人类偏好 → 训练 RM；再用 PPO + Value(critic) 稳定更新。
  \item \textbf{GRPO}：同题采样多答案 → 用规则 verifier 直接给分；用组内相对分数当 advantage → \textbf{不训练 RM，也不训练 Value}。
\end{itemize}
